% ----------------------------------------------------------
% Introdução
% Capítulo sem numeração, mas presente no Sumário
% ----------------------------------------------------------
%
% PENDÊNCIAS (apagar ao finalizar)
%  1. Declaração sobre o uso de ferramentas de IA: está comentada no fim
%     do arquivo e precisa ser revista antes de voltar ao texto.
%  2. Contribuição 6: confirmar com o capítulo de Resultados.
%  3. O capítulo de modelos de otimização não tem \label de capítulo,
%     por isso aparece por descrição em "Organização do trabalho".
%  4. Conferir se todas as chaves \cite existem no .bib.
% ----------------------------------------------------------

\chapter*[Introdução]{Introdução}
\addcontentsline{toc}{chapter}{Introdução}

A ciência da computação moderna tem suas bases nos trabalhos de Alan Turing na década de 1930, quando foi proposta uma noção abstrata de computação baseada em dispositivos programáveis, posteriormente denominada máquina de Turing \cite{sipser1996introduction, farias2013introduccao}. Esse modelo teórico estabelece que a computação pode ser descrita por meio da manipulação de símbolos em uma fita, de acordo com um conjunto finito de regras. Na prática, essa ideia fundamenta o funcionamento dos computadores digitais, nos quais a informação é representada em forma binária e processada por meio de instruções lógicas, sendo posteriormente interpretada em formatos compreensíveis ao usuário. \cite{hopcroftintroduction}

Entre a formalização teórica da computação e suas aplicações práticas, situam-se os chamados problemas combinatórios, que consistem na análise de um conjunto finito, porém frequentemente muito grande, de possíveis arranjos ou configurações, com o objetivo de identificar aquelas que satisfazem determinados critérios ou otimizam uma função específica \cite{papadimitriou1998combinatorial}. Esses problemas são caracterizados pelo crescimento exponencial do espaço de busca, tornando a exploração exaustiva e computacionalmente inviável. Nesse contexto, a ciência da computação busca desenvolver métodos e algoritmos capazes de lidar com essa complexidade, seja por meio de estratégias exatas, heurísticas ou aproximativas, estabelecendo uma ponte entre os modelos abstratos de computação e desafios concretos que demandam tomada de decisão eficiente.

Dentre as aplicações práticas desses modelos computacionais, a construção de rotas para entregas de mercadorias, o translado de passageiros, o transporte de cargas e o planejamento de trajetórias aéreas constituem problemas combinatórios de grande relevância econômica e operacional \cite{toth2002vehicle}. Tais problemas envolvem a busca por soluções ótimas ou quase ótimas em espaços de busca extensos, nos quais múltiplas variáveis e restrições devem ser consideradas simultaneamente, como custos, distâncias, tempo e capacidade dos veículos \cite{braekers2016vehicle}. A compreensão da dificuldade inerente a esses problemas, contudo, requer um arcabouço formal que permita classificá-los segundo sua complexidade computacional.

A teoria da complexidade computacional fornece um arcabouço formal para classificar a dificuldade desses problemas. Em particular, a classe NP (\textit{Nondeterministic Polynomial time}) engloba problemas de decisão para os quais uma solução candidata pode ser verificada em tempo polinomial, ainda que não se conheça, necessariamente, um método eficiente para encontrá-la. Já os problemas do tipo NP-difíceis correspondem àqueles que são, no mínimo, tão difíceis quanto os problemas mais complexos da classe NP, não havendo evidência de que possam ser resolvidos em tempo polinomial. Essa distinção é fundamental para compreender as limitações computacionais associadas a problemas de otimização combinatória e orientar o desenvolvimento de métodos de solução adequados. \cite{sipser1996introduction}

Inseridos nessa classe de problemas, os Problemas de Roteamento de Veículos (\textit{Vehicle Routing Problems} — VRP) destacam-se como uma classe fundamental de problemas de otimização combinatória \cite{psaraftis1995dynamic, laporte1992vehicle, Heidelberg2008travelingsalesmanperson}. Por serem classificados como NP-difíceis, não há evidência de que possam ser resolvidos em tempo polinomial à medida que o número de elementos envolvidos, tais como cidades, clientes, veículos, aumenta \cite{karp2009reducibility}. Essa dificuldade, contudo, não se estende de modo uniforme a todas as questões sobre uma solução candidata. Verificar se ela é viável, isto é, se satisfaz as restrições do problema, pode ser feito de forma eficiente \cite{sipser1996introduction, lenstra1981complexity}. Certificar que ela é ótima, por sua vez, não dispõe de método eficiente conhecido: exige estabelecer que nenhuma outra solução viável possui custo menor, o que reflete a dificuldade do próprio problema \cite{garey1990guide, papadimitriou1998combinatorial}. Além de seu interesse científico, associado ao desenvolvimento de métodos capazes de tratar problemas combinatórios de alta complexidade, os VRPs possuem forte relevância econômica e operacional, uma vez que soluções mais eficientes podem reduzir custos logísticos, melhorar a utilização de recursos e aumentar a qualidade do serviço prestado \cite{dantzig1959truck, toth2002vehicle}.

Um caso particular clássico do VRP é o Problema do Caixeiro Viajante (\textit{Traveling Salesman Problem} — TSP), que consiste em determinar a rota de menor custo para que um viajante percorra um conjunto de cidades, visitando cada uma exatamente uma vez e retornando à cidade de origem. Nesse problema, as distâncias entre as cidades são previamente definidas, e o objetivo principal é minimizar o custo total do percurso, frequentemente associado à distância total percorrida \cite{dantzig1954solution, applegate2011traveling}. Devido à sua natureza combinatória e à elevada complexidade associada à obtenção de soluções ótimas, o TSP constitui um importante caso de estudo para o desenvolvimento de novas abordagens computacionais.

Nesse contexto, a computação quântica surge como uma abordagem promissora para o tratamento de problemas combinatórios complexos \cite{albash2018adiabatic}. Enquanto os computadores clássicos representam a informação por meio de bits, nos quais valores discretos são associados fisicamente, por exemplo, à presença ou ausência de corrente elétrica em transistores \cite{shannon1948mathematical, sedra2011microelectronic}, nos computadores quânticos a informação é codificada em unidades denominadas qubits \cite{benioff1980computer, nielsen_quantum_2002}, cujo comportamento segue os princípios da mecânica quântica.

A teoria da computação quântica, cujas bases teóricas foram propostas por \textcite{benioff1980computer} e \textcite{feynman1982simulating}, fundamenta-se em princípios como superposição e entrelaçamento. A unidade fundamental e mais simples de informação nesse modelo é o \textit{qubit} (\textit{quantum bit}), que, diferentemente do bit clássico — restrito aos estados 0 ou 1 — pode ser descrito como uma combinação linear desses estados, permitindo representar simultaneamente múltiplas configurações \cite{nielsen_quantum_2002}.

Nesse contexto, surge a era dos dispositivos quânticos de escala intermediária e ruidosa (\textit{Noisy Intermediate-Scale Quantum} — NISQ), caracterizada por sistemas que operam com múltiplos qubits, ainda limitados tanto pela tecnologia de construção desses dispositivos quanto pela presença de ruído e pela ausência de correção completa de erros \cite{preskill2019quantum}. Esses dispositivos permitem a execução de algoritmos quânticos em pequena escala, embora com restrições quanto à profundidade dos circuitos e à precisão dos resultados \cite{bharti2022noisy, cheng2023noisy}.

Nesse cenário, diferentes paradigmas de computação quântica têm sido propostos e explorados, destacando-se o modelo de circuitos quânticos, baseado na aplicação sequencial de portas quânticas sobre qubits \cite{nielsen_quantum_2002}; a computação quântica adiabática, fundamentada na evolução contínua de sistemas quânticos para a obtenção de estados de menor energia \cite{albash2018adiabatic}; e a computação baseada em variáveis contínuas, na qual a informação é codificada em observáveis físicas contínuas, como posição e momento \cite{weedbrook2012gaussian}. Cada um desses paradigmas apresenta características distintas em termos de representação da informação, implementação física e estratégias algorítmicas, distinções que se tornam especialmente relevantes no contexto das limitações impostas pelos dispositivos NISQ, e que orientam diretamente o desenvolvimento de algoritmos quânticos com distintos requisitos de profundidade de circuito, tolerância a ruído e escalabilidade \cite{montanaro2016quantum}.

O aproveitamento dessas propriedades fundamenta o desenvolvimento de algoritmos quânticos capazes de tratar certos problemas de forma mais eficiente do que qualquer abordagem clássica conhecida. Entre os exemplos mais representativos dessa vantagem, destacam-se o algoritmo de Deutsch–Jozsa \cite{deutsch1991rapid}, que demonstra vantagem no modelo de consulta ao distinguir funções constantes de funções balanceadas com apenas uma única avaliação; o algoritmo de Grover \cite{grover1996fast}, que apresenta ganho quadrático em problemas de busca não estruturada, com complexidade $O(\sqrt{N})$; e o algoritmo de Shor \cite{shor1999polynomial}, que possibilita a fatoração de inteiros e o cálculo de logaritmos discretos em tempo polinomial no tamanho da entrada, ao passo que os melhores algoritmos conhecidos no paradigma computacional tradicional apresentam complexidade significativamente mais elevada, tornando esses problemas computacionalmente difíceis nesse contexto \cite{lenstra1993development}.

Uma estratégia recorrente para tratar problemas combinatórios em dispositivos quânticos consiste em codificar o problema em um Hamiltoniano cujo estado de menor energia corresponde à solução ótima. O exemplo canônico é a correspondência com o modelo de Ising, que permite escrever diversos problemas NP-difíceis como a minimização de uma energia sobre variáveis binárias \cite{lucas2014ising}, e que fundamenta tanto o recozimento quântico \cite{kadowaki1998quantum} quanto os algoritmos variacionais, entre os quais o \textit{Variational Quantum Eigensolver} (VQE) \cite{peruzzo2014variational} e o \textit{Quantum Approximate Optimization Algorithm} (QAOA) \cite{farhi_quantum_2014}, revistos em \cite{cerezo2021variational, moll2018quantum}. Os problemas de roteamento já foram tratados nesse arcabouço, com codificações do tipo QUBO para o TSP \cite{gonzalez2022gps} e para o VRP \cite{irie2019quantum, huang_solving_nodate}. Em todos esses casos, a codificação é construída sobre variáveis binárias.

O paradigma de variáveis contínuas oferece um substrato distinto para essa estratégia. Nele, a informação é codificada nas quadraturas de modos bosônicos, denominados qumodes, cujos operadores de posição $\hat{x}$ e de momento $\hat{p}$ possuem espectro contínuo e satisfazem a relação canônica $[\hat{x}, \hat{p}] = i\hbar$ \cite{weedbrook2012gaussian, braunstein2005quantum}. A computação sobre variáveis contínuas foi formalizada por \textcite{lloyd1999quantum}, e o formalismo é naturalmente implementado em plataformas fotônicas \cite{slussarenko2019photonic, fukui2022building}, para as quais há ferramentas de simulação estabelecidas \cite{killoran_strawberry_2019}. Entre os algoritmos de otimização propostos nesse paradigma, encontra-se uma versão do QAOA para problemas contínuos \cite{verdon2019quantum}. A representação de problemas combinatórios em qumodes é, por sua vez, menos explorada \cite{slysz2024solving, fingerhuth2018open}. Nela, a questão central é como representar escolhas discretas, como a posição de uma cidade na rota ou o veículo que a atende, em um sistema cujos observáveis são contínuos \cite{mcmahon2016fully}.

A investigação apresentada neste trabalho percorre esse problema em três passos, cada um motivado pela limitação estrutural do anterior. O primeiro consiste em verificar que a estratégia baseada em Hamiltonianos se transpõe ao paradigma de variáveis contínuas: a modelagem binária do TSP é transportada, termo a termo, para qumodes, associando um qumode a cada variável binária. A correspondência é viável, mas exige um número de qumodes que cresce com o quadrado do número de vértices e trata cada qumode como suporte de um único valor binário, o que desperdiça a riqueza de seu espaço de estados.

O segundo passo elimina essa redundância. Associa-se um qumode a cada posição da rota e a cidade é identificada pelo nível de ocupação na base de Fock, o que reduz o número de modos de $(N-1)^2$ para $N-1$ e produz um Hamiltoniano diagonal, cuja corretude pode ser demonstrada. Essa codificação, porém, emprega o qumode como um registrador de níveis discretos: a estrutura canônica do modo não desempenha papel algum, e a extensão a múltiplos veículos exigiria ampliar o rótulo discreto ao produto entre veículos e posições.

O terceiro passo supera essa limitação. Posição na rota e veículo são atribuídos às duas quadraturas de um mesmo qumode, por meio de variáveis modulares que, sob uma condição de compatibilidade entre escalas, tornam-se simultaneamente determináveis. Cada cidade passa a ocupar um qumode, e o número de modos permanece $N-1$ ainda que o problema representado, o VRP capacitado (\textit{Capacitated Vehicle Routing Problem} — CVRP), seja estritamente mais geral.

Dessa progressão decorrem as questões que orientam o trabalho. A primeira diz respeito à codificação: é possível construir, para problemas de roteamento, Hamiltonianos em variáveis contínuas que empreguem efetivamente as duas quadraturas do qumode e cujo estado fundamental, uma vez decodificado, corresponda a uma solução ótima? A segunda diz respeito ao algoritmo: um procedimento variacional de profundidade finita é capaz de preparar esse estado? A terceira diz respeito ao custo da representação, em número de modos e em dimensão de truncamento. As três são tratadas separadamente, pois uma resposta favorável à primeira não implica resposta favorável à segunda.

\section*{Objetivos}

O objetivo geral deste trabalho é propor e avaliar codificações Hamiltonianas de problemas de roteamento em variáveis contínuas, de modo que o estado fundamental corresponda a uma solução ótima e que a estrutura canônica do qumode seja efetivamente empregada, com ênfase no CVRP.

Os objetivos específicos são:
\begin{itemize}
    \item construir e analisar o mapeamento direto da modelagem binária do TSP para qumodes e a codificação posicional por níveis de Fock, demonstrando a corretude desta última;

    \item construir a codificação modular em duas quadraturas para o CVRP, estabelecendo as condições algébricas que a viabilizam e a regularização de energia finita que mantém a otimização em região fisicamente acessível;

    \item verificar a corretude das codificações por comparação entre o estado fundamental decodificado e o conjunto de soluções ótimas obtido por enumeração exaustiva clássica, em instâncias de pequeno porte;

    \item avaliar um algoritmo variacional, o VQE, sobre o Hamiltoniano modular, segundo dois critérios: o acerto combinatório e a proximidade do estado fundamental;

    \item caracterizar a dependência dos resultados em relação aos parâmetros da representação, a saber, a dimensão de truncamento, os pesos dos termos de penalização e a profundidade do circuito.
\end{itemize}

\section*{Contribuições}

As contribuições deste trabalho são as seguintes.
\begin{enumerate}
    \item \textbf{Redução do número de qumodes.} A codificação posicional por níveis de Fock representa um grafo de $N$ vértices com $N-1$ qumodes, contra $(N-1)^2$ do mapeamento direto da modelagem binária, e tem sua corretude demonstrada: para uma penalização $\lambda > N\,d_{\max}$, todo estado fundamental do Hamiltoniano corresponde a uma solução ótima do TSP.

    \item \textbf{Codificação modular em duas quadraturas.} Posição na rota e veículo são atribuídos às duas quadraturas de um mesmo qumode, por meio de variáveis modulares cuja comutação é garantida pela condição $ab = 2\pi\hbar$. O CVRP é representado com $N-1$ qumodes, o mesmo número da codificação anterior, e o TSP é recuperado como caso particular de um único veículo.

    \item \textbf{Hamiltoniano modular regularizado.} Construção dos termos de custo e de restrição (distância, colisão, lacuna, uso de veículos e capacidade), de um estabilizador que favorece a grade modular e de um termo de número de fótons que mantém a otimização em região de energia finita, com estimativa assintótica para a escolha de seu peso.

    \item \textbf{Tratamento da representação em dimensão finita.} Operadores de Weyl construídos a partir dos elementos de matriz exatos do operador de deslocamento, e ordenamentos simetrizados, em sanduíche para as penalidades e por anticomutador para os custos, que preservam hermiticidade e positividade sob truncagem.

    \item \textbf{Verificação da corretude por referência clássica independente.} O estado fundamental decodificado, obtido por diagonalização exata, é comparado com o conjunto ótimo da enumeração exaustiva clássica. Como essa enumeração opera diretamente sobre rotas, sem recorrer ao Hamiltoniano nem à regra de decodificação, a concordância nas instâncias testadas sustenta a corretude da codificação, e a discordância revelaria erro nela. Um avaliador de estados-produto, validado contra o cálculo direto, amplia a dimensão de truncamento acessível à análise.

    % TODO: confirmar este item com o capítulo de Resultados (dependência do corte).
    \item \textbf{Avaliação variacional em dois critérios.} {\color{red} verificar} O VQE é avaliado quanto ao acerto combinatório e quanto à proximidade do estado fundamental, o que permite separar limitações do algoritmo de limitações da codificação, e é caracterizada a dependência do desempenho em relação à dimensão de truncamento e aos pesos dos termos.
\end{enumerate}

Cabe delimitar o alcance do trabalho. Ele não tem como objetivo demonstrar vantagem computacional quântica sobre métodos clássicos. A compacidade das codificações propostas, medida em número de modos, não se confunde com aceleração: o espaço de configurações discretas permanece de cardinalidade exponencial no número de cidades, e o valor do Hamiltoniano de custo em uma configuração modular especificada pode ser avaliado classicamente. A contribuição visada é a proposição, a formalização e a avaliação de codificações em variáveis contínuas para problemas de roteamento, com análise estrutural comparativa em relação às codificações discretas existentes, o que contribui para a compreensão das vantagens e limitações de cada abordagem.

As avaliações são feitas por simulação em computador clássico, sobre instâncias de pequeno porte, limitadas a no máximo nove vértices. Esse limite decorre da capacidade computacional disponível, já que os recursos da simulação crescem com o número de vértices, e da estratégia de verificação: até essa ordem, a enumeração exaustiva fornece o conjunto completo das soluções ótimas, contra o qual as soluções recuperadas são comparadas. Para validar as codificações, essa dimensão é suficiente. A implementação em dispositivos físicos está fora do escopo.

Este trabalho está organizado como segue. O \autoref{sec:vrp} apresenta o VRP e sua modelagem matemática. O \autoref{sec:quantum_mechanics} reúne os fundamentos de mecânica quântica e de computação quântica, com ênfase no paradigma de variáveis contínuas. O capítulo seguinte discute modelos de otimização para problemas de rotas de veículos, em abordagens clássicas e quânticas. O \autoref{cap:formulacao} apresenta as três codificações desenvolvidas, em ordem de construção, cada uma encerrada pela análise estrutural que motivou a seguinte. O \autoref{cap:proposta} descreve os procedimentos empregados na avaliação, e o \autoref{cap:resultados} apresenta os resultados obtidos. O \autoref{cap:conclusao} reúne as conclusões e as perspectivas de continuidade.

A abordagem narrativa adotada neste trabalho privilegia a clareza e a progressividade na apresentação dos desenvolvimentos conceituais atuais e históricos desta área, para que seja um ponto de partida para aqueles que estejam iniciando neste estudo. Ao fazê-lo, visamos não apenas ampliar o escopo daqueles capazes de contribuir para esta área, mas também fornecer uma base sólida sobre a qual futuras pesquisas brasileiras possam ser construídas. Esta abordagem é fundamental não só para o enriquecimento do campo, mas também para a construção de uma comunidade acadêmica mais inclusiva e diversificada.

Durante a elaboração deste documento foram utilizadas ferramentas baseadas em Modelos de Linguagem de Grande Porte (\textit{Large Language Models} - LLMs) com o objetivo de auxiliar na revisão linguística, melhoria da fluidez textual e organização de trechos do texto. Essas ferramentas não foram utilizadas para geração de resultados científicos, desenvolvimento metodológico ou análise dos experimentos apresentados neste trabalho, sendo todo o conteúdo técnico e científico de responsabilidade dos autores.
